\chapter{Допуск family-to-multiplicity интертвейнера}
\label{ch:v8-baryon-c0-family-to-multiplicity-intertwiner-admission-gate}

Семейный qutrit, на котором действует $R_4^+$, является стандартным
вещественным представлением семейной группы, тогда как три connector-
направления текущей multiplicity-среды не несут объявленного семейного
действия:
\begin{equation}
 \mathcal E_{\rm fam}\simeq\mathbf3_{SO(3)_{\rm fam}},
 \qquad
 \mathcal E_{\rm mult}\simeq\mathbf1\oplus\mathbf1\oplus\mathbf1.
 \label{eq:v8-baryon-c0-family-multiplicity-representation-types}
\end{equation}
В стандартном базисе инфинитезимальные генераторы равны
\begin{equation}
 J_1=\begin{pmatrix}0&0&0\\0&0&-1\\0&1&0\end{pmatrix},\quad
 J_2=\begin{pmatrix}0&0&1\\0&0&0\\-1&0&0\end{pmatrix},\quad
 J_3=\begin{pmatrix}0&-1&0\\1&0&0\\0&0&0\end{pmatrix}.
 \label{eq:v8-baryon-c0-family-so3-generators}
\end{equation}
Для карты $T:\mathcal E_{\rm fam}\to\mathcal E_{\rm mult}$ условие
эквивариантности имеет вид
\begin{equation}
 TJ_a=0,\qquad a=1,2,3.
 \label{eq:v8-baryon-c0-family-multiplicity-intertwiner-equations}
\end{equation}
Точная линейная система содержит девять неизвестных и имеет
\begin{equation}
 \rank\mathcal C_{SO(3)}=9,\qquad
 \dim\ker\mathcal C_{SO(3)}=0,\qquad
 \operatorname{Hom}_{SO(3)}(\mathbf3,\mathbf1^{\oplus3})=0.
 \label{eq:v8-baryon-c0-family-multiplicity-so3-hom-zero}
\end{equation}
Следовательно, ненулевой перенос $R_4^+$ запрещён уже типами
представлений.

\section{Тетраэдрический остаток также не спасает карту}

Возьмём генераторы стандартной тройки $A_4$
\begin{equation}
 R=\operatorname{diag}(1,-1,-1),\qquad
 C=\begin{pmatrix}0&1&0\\0&0&1\\1&0&0\end{pmatrix},
 \qquad R^2=C^3=(RC)^3=I_3.
 \label{eq:v8-baryon-c0-family-a4-generators}
\end{equation}
Если target остаётся тривиальным, условия $TR=T$ и $TC=T$ дают
\begin{equation}
 \rank\mathcal C_{A_4}=9,\qquad
 \operatorname{Hom}_{A_4}(\mathbf3,\mathbf1^{\oplus3})=0.
 \label{eq:v8-baryon-c0-family-multiplicity-a4-hom-zero}
\end{equation}
Таким образом, даже замена непрерывной семейной группы её тетраэдрическим
остатком не создаёт допустимый ненулевой интертвейнер.

\section{Endpoint-метки и условное расширение}

Три текущих ветви различаются проекторами
\begin{equation}
 \Pi_0=\operatorname{diag}(1,0,0),\quad
 \Pi_1=\operatorname{diag}(0,1,0),\quad
 \Pi_2=\operatorname{diag}(0,0,1),\qquad
 \bigcap_i\Pi_i'=D_3,
 \label{eq:v8-baryon-c0-endpoint-projector-commutant}
\end{equation}
где $D_3$ --- диагональная алгебра. Непрерывное ортогональное
$SO(3)$-действие, сохраняющее каждый $\Pi_i$, должно лежать в дискретной
группе диагональных знаков, а потому на связной группе тождественно:
\begin{equation}
 \rho_{\rm mult}(SO(3))\subset D_3\cap O(3)
 =\{\operatorname{diag}(\pm1,\pm1,\pm1)\}
 \quad\Longrightarrow\quad \rho_{\rm mult}=I_3.
 \label{eq:v8-baryon-c0-endpoint-preserving-so3-triviality}
\end{equation}

Условный проход требует повысить multiplicity-среду до второй стандартной
тройки:
\begin{equation}
 \mathcal E_{\rm mult}^{\rm ext}\simeq\mathbf3_{SO(3)},
 \qquad
 T J_a=J_aT.
 \label{eq:v8-baryon-c0-promoted-multiplicity-so3-action}
\end{equation}
Тогда вещественная лемма Шура даёт
\begin{equation}
 \operatorname{Hom}_{SO(3)}(\mathbf3,\mathbf3)=\mathbb R I_3,
 \qquad T=\kappa I_3.
 \label{eq:v8-baryon-c0-promoted-intertwiner-schur-line}
\end{equation}
Совпадающие trace-метрики и изометрия сокращают эту прямую до
\begin{equation}
 T^{\mathsf T}T=I_3
 \quad\Longleftrightarrow\quad \kappa=\pm1,
 \qquad
 T R_4^+T^{\mathsf T}=R_4^+.
 \label{eq:v8-baryon-c0-promoted-isometric-intertwiner}
\end{equation}
Знак не влияет на перенесённый Hamiltonian. Однако само новое
$SO(3)$-действие смешивает endpoint-линии и несовместимо с сохранением
всех трёх $\Pi_i$ как неподвижных меток; оно является расширением parent,
а не следствием текущей алгебры.

Итоговый контракт равен
\begin{equation}
 \operatorname{Hom}_{G_{\rm current}}(\mathcal E_{\rm fam},
 \mathcal E_{\rm mult})=0,
 \qquad
 \dim\operatorname{Hom}_{G_{\rm extended}}=1,
 \qquad N_{\rm extension}=0/2.
 \label{eq:v8-baryon-c0-family-multiplicity-admission-ledger}
\end{equation}
Два необходимых новых объекта --- семейное действие на среде и его
совместимый подъём на endpoint-алгебру.

\begin{theorem}[Типовой запрет и условная линия Шура]
В текущем контракте не существует ненулевого $SO(3)$- или
$A_4$-эквивариантного отображения семейного qutrit в тривиальную
multiplicity-среду. После явного повышения target до стандартной тройки
пространство интертвейнеров становится одномерным и изометрия фиксирует
его с точностью до несущественного знака.
\end{theorem}

\begin{nogocriterion}[Нельзя забывать семейное действие]
Матрица $T\in M_3(\mathbb R)$ допустима только после удаления
$SO(3)_{\rm fam}$ из типа источника. Такое забывание превращает
представленческий запрет в произвольный девятипараметрический выбор и не
является выводом карты $c_0$.
\end{nogocriterion}

\begin{protocol}\begin{enumerate}
\item тип семейного qutrit: \textbf{$\mathbf3_{SO(3)}$};
\item текущий тип multiplicity-среды: \textbf{$\mathbf1^{\oplus3}$};
\item текущий $SO(3)$-Hom: \textbf{$0$};
\item текущий $A_4$-Hom: \textbf{$0$};
\item endpoint-preserving непрерывное действие: \textbf{тривиально};
\item условное target-расширение: \textbf{$\mathbf3_{SO(3)}$};
\item условный Hom после расширения: \textbf{$\mathbb RI_3$};
\item изометрический выбор: \textbf{$\pm I_3$};
\item новых структур расширения: \textbf{$2$};
\item выведено текущим parent: \textbf{$0/2$};
\item следующий гейт: \textbf{parent-origin семейного действия на среде}.
\end{enumerate}\end{protocol}

Машинный сертификат сохранён в
\path{s2t_v8_baryon_c0_family_to_multiplicity_intertwiner_admission_gate_results.json}.